% Part: first-order-logic
% Chapter: model-theory
% Section: substructures

\documentclass[../../../include/open-logic-section]{subfiles}

\begin{document}

\olfileid{mod}{bas}{sub}
\olsection{ذیلی \printtoken{p}{structure}}

اک !!a{structure}~$\Struct{M}$ دا !!{domain} دوجی ساخت~$\Struct{M'}$
دے دائرے دا ذیلی سیٹ ہو سکدا اے۔ پر $\Struct{M}$ نوں
$\Struct{M'}$ دا ”حصہ“ اُدوں ای سمجھنا چاہیدا اے جدوں نہ صرف
$\Domain{M}
\subseteq \Domain{M'}$ ہووے، بلکہ $\Struct{M}$ تے $\Struct{M'}$
زبان دیاں علامتاں دی تشریح وچ گھٹ توں گھٹ سانجھے حصے~$\Domain{M}$
اُتے اک دوجے نال ”متفق“ وی ہون۔
ایتھے زبان اوہی رہندی اے، پر دائرہ چھوٹا ہو سکدا اے؛
پچھلے گھٹائے روپ وچ زبان گھٹدی سی، دائرہ نہیں۔

\begin{defn}
\ollabel{defn:substructure}
!!{structure}s $\Struct M$ تے $\Struct M'$ اکّو زبان~$\Lang L$
لئی دتیاں ہون۔ اسیں $\Struct M$ نوں $\Struct M'$ دی
\emph{ذیلی !!{structure}} (substructure)، تے $\Struct M'$ نوں
$\Struct M$ دی \emph{توسیع} (extension) آکھدے آں، جس نوں
$\Struct M \substruct \Struct M'$ لکھدے آں، اُدوں تے صرف اُدوں جدوں
\begin{enumerate}
\item $\Domain{M} \subseteq \Domain{M'}$،
\item ہر مستقل $c \in \Lang L$ لئی، $\Assign{c}{M} =
    \Assign{c}{M'}$؛
\item ہر $n$-تھانواں والے !!{function} $f \in \Lang L$ لئی
  $\Assign{f}{M}(a_1, \dots, a_n) = \Assign{f}{M'}(a_1, \dots, a_n)$
  سارے $a_1$، \dots، $a_n \in \Domain{M}$ لئی۔
\item ہر $n$-تھانواں والے !!{predicate} $R \in \Lang L$ لئی، $\langle
  a_1, \dots, a_n\rangle \in \Assign{R}{M}$ اُدوں تے صرف اُدوں جدوں
  $\langle a_1, \dots,
  a_n\rangle \in \Assign{R}{M'}$ ہووے، سارے $a_1$، \dots، $a_n \in
  \Domain{M}$ لئی۔
\end{enumerate}
دائرہ غیر خالی اے، ہر مستقل دی قدر ایس وچ اے تے فنکشن دی چھوٹے
دائرے دیاں قدراں اُتے ورتوں دا نتیجہ وی ایتھے ای رہندا اے۔
ایہ گلاں دتی ہوئی ساخت دا دائرہ تے کل فنکشن ہون نال لازم نیں۔
\end{defn}

% PNB-SEM-035: پہلے درجے دی ساخت دا دائرہ غیر خالی اے؛خالی ذیلی سیٹ لئی ایہ دعویٰ نہیں۔
\begin{rem}
\ollabel{rem:substructure}
جے زبان وچ کوئی مستقل یا !!{function}s نہ ہون، تاں ہر غیر خالی $N
\subseteq \Domain{M}$ اک ذیلی !!{structure}~$\Struct{N}$
متعین کردا اے جو $\Struct M$ دی اے تے جس دا !!{domain}~$\Domain{N} = N$
اے؛ ایس لئی $\Assign{R}{N} =
\Assign{R}{M} \cap N^n$ رکھدے آں۔
\end{rem}

% ایس بارے کجھ ثابت کریے؟ مثالاں؟

\end{document}
